% RESEARCH_DRAFT_SHA256: 5333e79e937cee3a7a30c347506c2305f4dc429c57af95005a22c5347f7580a1
% Layout and source-reference wrappers expanded for the web renderer.
\documentclass[11pt]{article}
\usepackage[margin=25mm]{geometry}
\usepackage{amsthm}
\usepackage{fontspec,unicode-math}
\setmainfont{Latin Modern Roman}
\setmathfont{Latin Modern Math}
\usepackage{amsmath,amssymb,mathtools,xurl,hyperref}
\hypersetup{colorlinks=true,urlcolor=blue,linkcolor=blue}
\setcounter{secnumdepth}{0}
\setlength{\parindent}{0pt}
\setlength{\parskip}{5pt}
\setlength{\emergencystretch}{3em}
\newcommand{\sref}[2]{\hyperlink{source:#1:#2}{[S#2]}}
\newcommand{\cataloguescope}[1]{\paragraph{Source record.}#1}
\newcommand{\floor}[1]{\left\lfloor#1\right\rfloor}
\newtheorem{proposition}{Proposition}
\newtheorem{lemma}[proposition]{Lemma}
\begin{document}
\paragraph{Research manuscript.} OpenAI \texttt{gpt-6-astra}, reasoning \texttt{ultra}; 11 October 2026.
\paragraph{Preserved source body.} The complete catalogue mathematical text and sources follow with layout and source links expanded, including source qualifications and its historical preparation notes. New research begins at the section headed Research attempt.

% UnsolvedMath ID 30002240; source catalogue code OWR-12177-009
\section{Unique Games Conjecture}\label{30002240}
{\small UnsolvedMath ID 30002240 · Computer Science}
\subsection{Definitions and mathematical statement}

\paragraph{Shared notation.}Standard number systems are $\mathbb N,\mathbb Z,\mathbb Q,\mathbb R,\mathbb C$. Specialized notation and index ranges follow the cited source and the definitions below.


A Unique Games instance assigns labels from a finite alphabet to graph vertices, with each edge imposing a permutation relation on the endpoint labels. Its value is the largest fraction of satisfiable edges. The standard conjecture asks whether, for every positive tolerance, a sufficiently large alphabet makes distinguishing near-unit value from near-zero value NP-hard.
\paragraph{Record scope and qualifications.}
The catalogue title is retained as metadata. Named certificates, contracts, bounds or frameworks in that title are not asserted as results of the original source extract.
This repeats the Unique Games target already represented by canonical source records 43 and 30001374.

Source-record limitation: Duplicate Unique Games target, already present under IDs 43 and 30001374.

Related catalogue records: \href{https://mehmetmars7.github.io/Erdosproblems-llm-hunter/problem.html?type=open_problems&id=43}{43}, \href{https://mehmetmars7.github.io/Erdosproblems-llm-hunter/problem.html?type=open_problems&id=30001374}{30001374}. Separate source IDs are preserved for their provenance.
\paragraph{Statement recorded in the supplied source.}
Is Khot's Unique Games Conjecture true? \href{https://doi.org/10.4171/owr/2012/54}{[S2]}
\paragraph{Source status.}
The supplied catalogue labels this source record open; that metadata was checked against the supplied snapshot on 2026-10-11. The label is preserved for provenance and is not an independent verification of present open status or mathematical correctness. The catalogue literature assessment is dated 2026-08-21.
\subsection{Short English statement}
Ask whether nearly satisfiable permutation-constraint instances are hard to distinguish from ones with very low satisfiable value.
\subsection{Sources}
\begin{itemize}
\item[S1] Ulam AI and the UnsolvedMath Contributors, UnsolvedMath: A Curated Collection of Open Mathematics Problems, record 30002240 (OWR-12177-009). Catalogue curation and metadata: CC BY 4.0. The source extract is retained; definitions, English summaries and qualifications are editorial additions.
\url{https://huggingface.co/datasets/ulamai/UnsolvedMath}.
\item[S2] Complexity Theory (2013), Oberwolfach Reports, DOI 10.4171/owr/2012/54, pp. 3267--3304
\url{https://doi.org/10.4171/owr/2012/54}.
\item[S3] Khot original formulation
\url{https://doi.org/10.1145/509907.510017}.
\item[S4] current status survey
\url{https://theoryofcomputing.org/articles/v021a010/v021a010.pdf}.
\item Source-specific Oberwolfach Reports licences, including CC BY-SA 4.0 where applicable, are retained; see THIRD\_PARTY\_NOTICES.md.
\end{itemize}
\paragraph{Preparation and review.}
The original source and UnsolvedMath contributors retain attribution for the source extract. Codex (GPT-6; exact serving revision not exposed) assisted with definitions, English summaries, scope notes and formatting. No new human mathematical authorship or independently verified solution is asserted. Permanent source ID 30002240 and code OWR-12177-009 checked against the complete supplied registry and current repository. Distinct upstream records, including repeated questions and broad or damaged extracts, are retained for complete coverage of the source CS category. A repeated record is not a new independent conjecture. Complete mathematical-equivalence and current-literature review has not been performed. Related canonical identities: 43, 30001374. Preparation and final review follow contribution prompts C and D, with source gaps disclosed. Structural and rendering checks do not establish a new theorem or open status.
\label{end:30002240}

\section{Research attempt}
\subsection{Literature, source context and scope}
The full original catalogue body and source qualifications are preserved above. The relevant primary workshop passages were read in the cached publisher report. Historical claims and subsequent attributed results are distinguished from the mathematical work below. A restricted example, standard lemma or known theorem is not counted as new progress on the entire original target. Missing assumptions and repeated or rhetorical extracts are not silently turned into different solved problems.

\subsection{Exact consistency does not solve the near-perfect gap}
The source surveys the Unique Games Conjecture and semidefinite approaches. We test an elementary propagation route. Consider a finite graph with alphabet $[q]$, $q\ge1$, and each oriented edge $uv$ carrying a permutation $\pi_{uv}$, with the reverse constraint its inverse. A labeling satisfies this edge when $\ell(v)=\pi_{uv}(\ell(u))$.
\begin{proposition}
Whether all edges can be satisfied is decidable in $O(q(|V|+|E|))$ permutation evaluations. However, this exact test rejects instances whose optimum satisfied fraction tends to one.
\end{proposition}
\begin{proof}
In each connected component choose a root and a spanning tree. For every possible root label, propagate labels along tree edges using their permutations or inverses, then check every edge. Every fully satisfying labeling is uniquely determined by its root label along that tree, so the procedure succeeds exactly when a fully satisfying labeling exists. The claimed cost follows by trying $q$ labels in each component and summing their sizes.
For the second assertion take a cycle of length $L\ge3$ with alphabet $\{0,1\}$. Give $L-1$ edges identity constraints and the remaining edge the flip $b\mapsto1-b$. Satisfying all edges would propagate a label around the cycle to its complement, an impossibility. The constant-zero labeling satisfies all identity edges, so the optimum is exactly $(L-1)/L$.
\end{proof}
The attempted inference ``efficient exact consistency testing refutes the conjecture'' therefore fails. For every fixed $\eta>0$, a cycle with $L>1/\eta$ is a near-perfect instance, yet the exact test rejects it. A promise distinction between value at least $1-\eta$ and value at most a small constant is a different task; the cycle example alone gives neither an algorithm nor a hardness proof for that distinction.
\subsection{Verification and remaining gaps}
All binary labelings of the frustrated cycles with $3\le L\le10$ were enumerated, obtaining optimum $L-1$ satisfied edges. The propagation algorithm was also checked against exhaustive satisfiability for these examples. No new gap algorithm, semidefinite hierarchy analysis, or complexity lower bound is established. The 2012 survey's historical algorithmic claims are not presented as current exhaustive knowledge.

\section{References and source attribution}
Primary workshop report: \url{https://doi.org/10.4171/owr/2012/54}. Relevant passages read: 2012/54 p. 3286, Boaz Barak’s survey of Unique Games, subexponential algorithms and sum-of-squares candidates.

All retrievals were checked on 11 October 2026. Original question attribution remains in the preserved Sources section. External theorems are attributed and not claimed as new discoveries or fully verified merely by citation.


\paragraph{Integration review.} Codex (GPT-6; exact serving revision not exposed) reviewed the written arguments and expanded layout and source-reference wrappers. Repository metadata and the production web renderer were checked using Python 3.13 and JavaScript with MathJax 3.2.2. This model-assisted review is not independent human proof verification; no PDF compilation or formal proof checking is asserted.

\section{Completion Estimate}
\noindent\textbf{COMPLETION: 0\%}
No substantiated new progress toward the entire original target is claimed. The worked results and failed route are delimited above; a known result, illustrative model, or auxiliary obstruction is not counted as original-target completion.

\section{Final status and attribution}
\textbf{Attempt status: unresolved.} This describes this research attempt, not a new ruling on the source catalogue's status or on external literature claims.
The mathematics and this self-audit were generated by OpenAI \texttt{gpt-6-astra}, reasoning \texttt{ultra}, on 11 October 2026. Codex (GPT-6; exact serving revision not exposed) orchestrated the batch, reviewed the written mathematical arguments, and contributed formatting/repository review. No human mathematical authorship, independent human proof checking, or formal verification is claimed. The historical preparation paragraphs in the reproduced catalogue body concern that body only.

Contribution prompt A was applied to the complete definition and supplied source record, followed by the adversarial checks of prompt B. Local reproduction files are \texttt{batch019/checks.py}, \texttt{computations.json}, and the per-ID audit JSON. They were not public links at preparation time. Numerical checks are finite and do not prove unrestricted results. No LaTeX engine was found locally; compilation was not performed by this research agent.

\end{document}
